\documentclass{article}
\usepackage{geometry}
\usepackage{graphicx} % Required for inserting images
\usepackage{amsmath, amsthm, amssymb}
\usepackage{parskip}
\newgeometry{vmargin={15mm}, hmargin={24mm,34mm}}
\theoremstyle{definition} 
\newtheorem{definition}{Definition}

\newtheorem{theorem}{Theorem}[section]
\newtheorem{lemma}[theorem]{Lemma}
\newtheorem{proposition}[theorem]{Proposition}
\newtheorem{example}[theorem]{Example}
\newtheorem{corollary}{Corollary}[theorem]
\newcommand{\nil}{\operatorname{nil}}
\newcommand{\ann}{\operatorname{Ann}}
\newcommand{\codim}{\operatorname{codim}}

\title{Dimension in Commutative Algebra}
\date{April 2025}
\author{Boran Erol}

\begin{document}

\maketitle

% Borcherds has a lecture where he gives 100 definitions of dimension and shows why they all suck.

\begin{lemma}
    $ \codim(P) = \dim(R_{P}) $.
\end{lemma}
\begin{proof}
    Prime ideals of $ R $ contained in $ P $ are in inclusion-respecting
    bijective correspondence with prime ideals of $ R_{P} $, so we have
    the desired result.
\end{proof}

\begin{lemma}
    The dimension of $ R $ is the supremum of
    the dimensions of $ R_{P} $ where $ P $
    is a maximal ideal of $ R $.

    In other words, dimension is a local concept.
\end{lemma}

\begin{lemma}[Gathmann CA Lemma 11.6]
    Let $ R $ be a ring of finite dimension in which all maximal chains
    of prime ideals have the same length. Let $ P $ be a prime ideal of $ R $.

    Then, $ R/P $ is a ring of finite dimension where all maximal chains of
    prime ideals have the same length.

    In particular, we have that $ \dim(R) = \dim(R/P) + \dim(R_{P}) $.
\end{lemma}

\begin{corollary}
    $ \dim(R_{P}) = \dim(R) $ if and only if $ P $ is a maximal ideal.
\end{corollary}

Let's first show that integral ring extensions preserve the dimension of a ring.

\begin{proposition}[Gathmann CA Proposition 11.9]
    All maximal chains of prime ideals in $ K[x_{1}, \ldots, x_{n}] $
    have length $ n $.
\end{proposition}
\begin{proof}
    For $ n = 0 $, the statement is trivial.
    For $ n = 1 $, the statement follows from the fact that $ K[x] $ is a PID.


\end{proof}

\begin{corollary}
    Let $ R $ be a finitely generated algebra over a field $ K $.
    Assume $ R $ is an integral domain.
    Then, every maximal chain of prime ideals in $ R $ has length $ \dim(R) $.
\end{corollary}
\begin{proof}
    We can write $ R $ as $ K[x_{1}, \ldots, x_{n}] / P $ for some $ n \geq 1 $
    and prime ideal $ P $ of $ R $.   % TODO: Finish the proof yourself, I read and understood it.  
\end{proof}

Let's now convert this statement to geometry.

Let $ Y $ be an irreducible subvariety of an irreducible variety $ X $.
Then, $ \dim(X) = \dim(A(X)) = \dim(A(X)/I(Y)) + \codim(I(Y)) = \dim(A(Y)) + \codim_{X}(Y) = \dim(Y) + \codim_{X}(Y)$.

Moreover, the local dimension of $ X $ is $ \dim(X) $ at every point of $ X $. In other words,
all localizations of $ A(X) $ at maximal ideals have dimension $ \dim(X) $.

\begin{definition}
    Let $ R $ be a ring. Let $ P $ be a prime ideal of $ R $ and $ n \geq 1 $.
    Let $ P^{(n)} := \{ a \in R : ab \in P^{n} \text{ for some $ b \in R/P $} \} $.
    $ P^{(n)} $ is called the \textbf{$ n $-th symbolic power of $ P $}.
\end{definition}



\end{document}
